\question[8]
\lstinline|heap-three.c| (es.\ 19) fa tre \lstinline|malloc(32)| (\lstinline|a|, \lstinline|b|, \lstinline|c|), tre \lstinline|strcpy| senza limiti dai primi tre argomenti, poi \lstinline|free(c); free(b); free(a)| e infine una \lstinline|printf| (tradotta dal compilatore in una \lstinline|call| a \lstinline|puts@plt|). \lstinline|winner()| esiste ma non è chiamata.
(a) Perché va linkato contro dlmalloc 2.7.2 e solo i386? (b) Qual è il bersaglio della \emph{write-what-where} via \lstinline|unlink()| e come si arriva a \lstinline|winner()|? (c) Perché una \lstinline|size| fittizia ``da manuale'' (64) non basta a innescare \lstinline|unlink()|?
\begin{solution}
(a) La glibc di bookworm (2.36) ha il \emph{safe unlinking}: l'attacco classico sarebbe bloccato. Il commento originale di Phoenix dichiara il livello ``linked against malloc-2.7.2.c'' (dlmalloc di Doug Lea, SHA1 \lstinline|407329d1...|): \lstinline|build.sh| compila quella vecchia \lstinline|malloc| come oggetto e la linka al binario (per interposizione \lstinline|malloc|/\lstinline|free| sono \emph{definiti} nel binario, non importati da libc), e la sua \lstinline|unlink()| non ha controlli di integrità. Solo i386 (\lstinline|gcc -m32|): a 64 bit l'attacco richiederebbe una strategia diversa, non trattata.\\
(b) Bersaglio: la entry \textbf{GOT di \texttt{puts()}} (\lstinline|0x0804e018|, da \lstinline|objdump -R|). Sovrascrivendo i metadati del chunk ``successivo'' (\lstinline|size| con \lstinline|PREV_INUSE| a 0, \lstinline|fd|/\lstinline|bk| falsi), \lstinline|free()| crede quel chunk libero e invoca \lstinline|unlink()|, che scrive nella GOT di \lstinline|puts| l'indirizzo di un \emph{trampolino} (\lstinline|push winner; ret|, posto sullo stack dentro \lstinline|argv[1]|). La \lstinline|printf| finale, eseguita come \lstinline|puts@plt|, salta nel trampolino e di lì in \lstinline|winner()| (\lstinline|0x08049216|). Non si punta \lstinline|bk| direttamente a \lstinline|winner| perché \lstinline|unlink| ne corromperebbe un campo.\\
(c) Perché una \lstinline|size| sotto la vera soglia \lstinline|max_fast| manda il chunk in un \emph{fast bin} (riuso semplice, nessun consolidamento) invece che sul percorso che chiama \lstinline|unlink()|. Su questo binario la soglia reale, letta a runtime da \lstinline|av_.max_fast|, è \textbf{72 byte} (non i 64 ``da manuale'', per via del \lstinline|request2size()| interno di dlmalloc): la \lstinline|size| fittizia va scelta sopra 72 e con \lstinline|PREV_INUSE| coerente.
\end{solution}
